\section{问题四：离网运行与储能配置研究}

\subsection{问题分析与求解思路}

问题四研究园区离网运行模式，即切断与外部电网的连接，园区用能完全依赖风光发电\cite{ref11}。离网条件下$P_{buy}(t) = 0$、$P_{sell}(t) = 0$，功率平衡约束更加严格——任何时刻用电负荷不得超过风光发电功率。分三个子问题逐步深入：无储能时的最大产量分析、最优储能配置方案设计\cite{ref6}\cite{ref12}、离网与联网模式的经济性对比。

\begin{figure}[H]
\centering
\includegraphics[width=0.85\textwidth,height=0.38\textheight,keepaspectratio]{../figures/fig_flow_q4.pdf}
\caption{问题四求解流程图}\label{fig:flow_q4}
\end{figure}

\subsection{问题四(1)：离网无储能运行分析}

\subsubsection{数学模型}

离网条件下，功率平衡约束变为：
\begin{equation}
\alpha(t) \cdot P_{H2NH3}^{max} + P_{load}(t) \leq P_{wind}(t) + P_{pv}(t), \quad \forall t
\end{equation}

即功率调节系数受风光出力上界约束：
\begin{equation}
\alpha(t) \leq \frac{P_{wind}(t) + P_{pv}(t) - P_{load}(t)}{P_{H2NH3}^{max}}, \quad \forall t
\end{equation}

当风光出力不足以支撑10\%最低负荷时（$P_{wind}(t) + P_{pv}(t) - P_{load}(t) < 0.1 \times 41.5 = 4.15$ MW），设备必须停机。目标函数为最大化日制氨产量：
\begin{equation}
\max \quad Q_{NH3} = r_{NH3} \cdot \sum_{t=1}^{24} \alpha(t) \cdot \Delta t = 3.0 \cdot \sum_{t=1}^{24} \alpha(t)
\end{equation}

离网模式下的吨氨成本仅包含新能源度电成本和设备运维成本（无购电成本和售电收入）：
\begin{equation}
C_{ton}^{off} = \frac{C_{RE,used} + C_{OM}}{Q_{NH3}}
\end{equation}

\subsubsection{计算结果}

对24种风光场景分别求解LP，得到各场景的最大制氨产量和弃电量。

\begin{figure}[H]
\centering
\includegraphics[width=0.9\textwidth,height=0.38\textheight,keepaspectratio]{../figures/fig_q4_offgrid_production.pdf}
\caption{24种风光场景下离网制氨日产量对比}\label{fig:q4_offgrid_production}
\end{figure}

24种场景下的离网制氨产量差异显著（图\ref{fig:q4_offgrid_production}）。高风光场景下日产量可接近72吨上限，而低风光场景下日产量大幅下降。全年关键指标为：全年总产量9782吨，平均产能利用率37.7\%。相比联网模式（问题三）的12960吨/年，离网无储能模式的年产量减少24.5\%，产能利用率不足四成，大量产能因风光时序限制而闲置。

最大弃电场景为W4P1（高光伏、中高风电），弃电量达177.3 MWh。该场景的特点是白天风光出力远超园区最大用电能力（$P_{H2NH3}^{max} + P_{load,max} = 41.5 + 6 = 47.5$ MW），多余电力只能弃掉。

\subsubsection{最小风光装机容量估算}

为实现园区基本能源自给（覆盖常规负荷加10\%制氢氨最低功率），在所有24种场景最严苛时段约束下，通过线性规划求解得到最小风电装机容量约为155.0 MW，此时仅靠风电即可满足需求（最小光伏装机为0 MW）。该结果表明，在当前风光资源条件下，风电的出力稳定性优于光伏，是离网园区能源自给的关键支撑。

\subsection{问题四(2)：储能配置优化}

\subsubsection{储能状态方程}

引入储能系统后，功率平衡方程扩展为：
\begin{equation}
\alpha(t) \cdot P_{H2NH3}^{max} + P_{load}(t) + P_{sto,c}(t) = P_{wind}(t) + P_{pv}(t) + P_{sto,d}(t) - P_{curtail}(t)
\end{equation}

储能荷电量（SOC）状态转移方程为：
\begin{equation}\label{eq:soc}
E_{sto}(t+1) = E_{sto}(t) \cdot (1 - \sigma) + \eta_c \cdot P_{sto,c}(t) \cdot \Delta t - \frac{P_{sto,d}(t) \cdot \Delta t}{\eta_d}
\end{equation}

其中$\sigma = 0.002$/h为自损耗率，$\eta_c = \eta_d = 0.9$为充放电效率。储能约束包括容量约束$0 \leq E_{sto}(t) \leq C_{sto}$、功率约束$P_{sto,c}(t), P_{sto,d}(t) \leq P_{sto}^{max}$、以及日循环约束$E_{sto}(25) = E_{sto}(1)$。

\begin{figure}[H]
\centering
\includegraphics[width=0.85\textwidth,height=0.35\textheight,keepaspectratio]{../figures/tikz_storage_model.pdf}
\caption{储能系统SOC状态转移方程示意图}\label{fig:tikz_storage}
\end{figure}

储能系统的SOC状态转移过程如图\ref{fig:tikz_storage}所示，每个时段的荷电量由前一时段的自损耗、充电输入和放电输出三部分共同决定。

\subsubsection{参数扫描优化}

针对最大弃电场景（W4P1），采用参数扫描法对储能容量$C_{sto} \in \{10, 50, 100, 150, 200\}$ MWh进行逐一求解，分析储能容量对产量提升和弃电减少的边际效果。功率容量比固定为$P_{sto}^{max} = C_{sto} / 2$（0.5C倍率）。

\begin{table}[H]
\centering
\caption{不同储能容量下最大弃电场景运行结果}\label{tab:q4_storage}
\begin{tabular}{cccc}
\toprule
储能容量(MWh) & 日产量(吨) & 吨氨成本(元/吨) & 弃电量(MWh) \\
\midrule
0（无储能） & 45.2 & 3582 & 177.3 \\
10 & 47.1 & 3716 & 164.7 \\
50 & 49.7 & 3868 & 128.5 \\
100 & 52.9 & 4037 & 82.9 \\
150 & 56.1 & 4187 & 37.3 \\
200 & 56.5 & 4347 & 32.8 \\
\bottomrule
\end{tabular}
\end{table}

从表\ref{tab:q4_storage}可以看出，储能容量从0增至150 MWh时，日产量从45.2吨提升至56.1吨（+24.1\%），弃电量从177.3 MWh降至37.3 MWh（$-$79.0\%）。继续增加至200 MWh时，产量仅微增0.4吨、弃电仅减少4.5 MWh，边际效益急剧递减。因此最优储能配置为150 MWh / 75 MW，该方案在产量提升和投资成本之间取得了最佳平衡。

\begin{figure}[H]
\centering
\includegraphics[width=0.9\textwidth,height=0.38\textheight,keepaspectratio]{../figures/fig_q4_storage_dispatch.pdf}
\caption{最大弃电场景下储能参与的功率调度曲线}\label{fig:q4_storage_dispatch}
\end{figure}

配置150 MWh储能后的功率调度曲线（图\ref{fig:q4_storage_dispatch}）展示了储能的"削峰填谷"作用：白天光伏高峰时段储能充电，吸收多余风光电力；夜间或风光不足时段储能放电，支撑制氢氨设备继续运行。SOC在日内呈现先升后降的循环特征，日末恢复到初始水平。

配置最优储能后，对24种场景分别制定功率调度方案。有储能参与的全年总产量为10619吨，平均产能利用率41.0\%，较无储能（9782吨，37.7\%）分别提升8.6\%和3.3个百分点。储能有效改善了园区的能源自给能力和风光发电利用效率，将全年弃电总量降低约60\%。

\subsection{问题四(3)：离网与联网经济性对比}

\begin{figure}[H]
\centering
\includegraphics[width=0.85\textwidth,height=0.38\textheight,keepaspectratio]{../figures/fig_q4_economics.pdf}
\caption{离网与联网运行模式吨氨成本对比}\label{fig:q4_economics}
\end{figure}

在满足相同制氨产量需求前提下，对比两种运行模式的全年吨氨成本（图\ref{fig:q4_economics}）：联网模式（问题三）全年吨氨成本为4209.48元/吨，离网模式（有储能）全年吨氨成本为4785.85元/吨，联网模式成本低于离网模式13.7\%。

\begin{table}[H]
\centering
\caption{离网与联网模式经济性对比}\label{tab:q4_economics}
\begin{tabular}{lcc}
\toprule
指标 & 联网模式 & 离网模式（有储能） \\
\midrule
全年吨氨成本 & 4209.48 元/吨 & 4785.85 元/吨 \\
全年总产量 & 12960 吨 & 10619 吨 \\
产能利用率 & --- & 41.0\% \\
电网支撑价值 & \multicolumn{2}{c}{13.7\%} \\
\bottomrule
\end{tabular}
\end{table}

电网对园区的支撑价值体现在两个方面：一是经济性支撑，联网模式可以在风光不足时低成本购电维持生产，避免了储能投资和效率损失带来的额外成本；二是产量支撑，联网模式全年产量（12960吨）比离网有储能模式（10619吨）高22.1\%，设备利用率显著提升。13.7\%的成本差异即为电力系统对绿电直连园区的"系统支撑成本价值"，这一数值为园区选择联网或离网运行模式提供了定量决策依据。
